\documentclass[11pt]{article}

\usepackage[a4paper,margin=1in]{geometry}
\usepackage{amsmath,amssymb,amsfonts,amsthm,mathtools,bm}
\usepackage{booktabs}
\usepackage{multirow}
\usepackage{array}
\usepackage{enumitem}
\usepackage{xcolor}
\usepackage{hyperref}
\usepackage{graphicx}
\usepackage{float}
\usepackage{setspace}
\usepackage{titlesec}

\hypersetup{
    colorlinks=true,
    linkcolor=blue,
    citecolor=blue,
    urlcolor=blue
}

\setstretch{1.08}
\setlength{\parskip}{0.45em}
\setlength{\parindent}{0pt}

\titleformat{\section}{\large\bfseries}{\thesection.}{0.5em}{}
\titleformat{\subsection}{\normalsize\bfseries}{\thesubsection.}{0.5em}{}
\titleformat{\subsubsection}{\normalsize\bfseries}{\thesubsubsection.}{0.5em}{}

\title{\textbf{LoTR: Logic-of-Thought Routing for Plug-in Reasoning}\\
\vspace{0.3em}
\large A Detailed Proposal}
\author{}
\date{}

\begin{document}

\maketitle

\section{Project Positioning}

This project studies a simple question: \emph{can reasoning be improved by dynamically matching the model's current internal state with a more suitable inference pathway at each step?}

Existing reasoning methods are often built around \emph{external program injection}: they specify a reasoning scaffold from outside the model, such as chain-of-thought prompting, self-consistency sampling, self-refinement loops, or tree/search-based exploration. These methods differ in how they organize trajectories, but they usually do not explicitly model the model's own \emph{internal state evolution} during reasoning, nor do they adapt the model's active inference pathway according to the current step-level logic demand.

Our central view is that reasoning quality depends not only on the external scaffold, but also on whether the model is using a suitable \emph{internal pathway} for the current step. At different reasoning steps, the model may be in different latent logic states. Some attention heads may support the current logic, some may be irrelevant, and some may even interfere with the current reasoning move. Therefore, a fixed pathway is suboptimal. We should instead infer the current logic state from internal information transfer, and softly route the model toward a more suitable pathway for that step.

Based on this intuition, we propose \textbf{Logic-of-Thought Routing (LoTR)}, a \emph{plug-in reasoning controller} that can be attached to diverse reasoning paradigms. LoTR does not replace CoT, SC, SR, MCTS, or other reasoning frameworks. Instead, it acts as a lightweight internal routing layer: at each sentence-level reasoning step, it reads the model's current internal state, predicts a distribution over a compact set of global logic states, and then softly combines precomputed logic-specific routing templates to modulate attention heads for the next step.

The method is intentionally simple. It has two main ingredients:
\begin{enumerate}[leftmargin=1.5em]
    \item a lightweight \emph{logic probe} that maps the current reasoning state to a probability distribution over logic states;
    \item a cached set of \emph{logic-specific routing templates} that determine how strongly each attention head should participate under each logic state.
\end{enumerate}

The resulting method is plug-in, lightweight, reusable, and compatible with multiple reasoning paradigms.

\section{Core Thesis}

The core thesis of this project is the following:

\begin{quote}
\textbf{Reasoning should be supported by step-adaptive internal routing.} A model's current reasoning step can be associated with a latent logic state, and different logic states require different internal attention-head participation patterns. If we can identify a compact set of global logic states and their corresponding routing templates, then we can improve reasoning by dynamically selecting a more suitable pathway at each step.
\end{quote}

This thesis can be decomposed into three claims:

\begin{enumerate}[leftmargin=1.5em]
    \item \textbf{Global logic states exist.} Across diverse problems and reasoning paradigms, there exists a compact set of recurring internal logic regimes that can be discovered from reasoning trajectories.
    \item \textbf{Logic states are coupled with sparse routing patterns.} Different logic states are associated with different attention-head participation profiles, while many heads remain logic-insensitive.
    \item \textbf{Logic-conditioned routing improves plug-in reasoning.} Softly routing attention heads according to the current logic state improves reasoning quality more generally than static routing, fixed pruning, or external scaffolds alone.
\end{enumerate}

\section{Design Principles}

The method should remain clean and paradigm-like. We therefore impose the following design principles.

\subsection{Plug-in Rather than Replacement}

LoTR must be attachable to existing reasoning paradigms without redefining them. If a base method uses CoT, SC, SR, or MCTS, LoTR should simply operate on top of its stepwise generation process.

\subsection{Internal-State-Driven Rather than Rule-Driven}

We do not define logic states manually by symbolic rules, nor do we hard-code specific logic programs. Logic states are discovered from internal reasoning trajectories.

\subsection{Finite Logic Basis Rather than Instance-Level Optimization}

We do not optimize a new controller for each problem instance. Instead, we discover a compact global set of logic states offline, compute reusable routing templates for them, and perform only lightweight online combination at inference time.

\subsection{Soft Routing Rather than Hard Switching}

Different reasoning steps may exhibit mixed logic demand. Therefore, the online controller should output a probability distribution over logic states and softly combine their routing templates.

\subsection{Minimal Controller Complexity}

The online controller should be simple:
\begin{itemize}[leftmargin=1.5em]
    \item no extra search,
    \item no extra planner,
    \item no separate memory module,
    \item no handcrafted fallback branch,
    \item no multi-stage controller stack.
\end{itemize}
The core online computation should ideally reduce to one probe evaluation and one weighted combination over cached templates.

\section{Problem Formulation}

Let $f_\theta$ be a frozen autoregressive language model with $L$ layers and head set
\[
\mathcal{H} = \{(\ell,h)\mid \ell=1,\dots,L,\ h=1,\dots,H_\ell\}.
\]

Consider a reasoning paradigm $\Pi$ (for example, CoT, SC, SR, or MCTS). For an input problem $x$, $\Pi$ induces a stepwise reasoning trajectory
\[
y_{1:T} = (y_1, y_2, \dots, y_T),
\]
where each $y_t$ is a sentence-level reasoning unit.

Our goal is not to redesign $\Pi$, but to augment it with a plug-in controller that modifies the model's internal routing at each step. Specifically, for each step $t$, we aim to infer a logic-state distribution
\[
p_t \in \Delta^{K-1},
\]
where $K$ is the number of global logic states, and use it to construct a routing matrix
\[
W_t \in [0,1]^{L\times H},
\]
which determines the degree to which each head participates in the next-step computation.

The resulting controller should satisfy:
\begin{enumerate}[leftmargin=1.5em]
    \item it operates at sentence level;
    \item it is independent of the outer reasoning scaffold;
    \item it reuses a fixed set of cached logic templates;
    \item it only changes attention-head participation weights.
\end{enumerate}

\section{Method Overview}

LoTR consists of two stages:

\begin{enumerate}[leftmargin=1.5em]
    \item \textbf{Offline Logic--Pathway Compilation}
    \begin{itemize}[leftmargin=1.5em]
        \item collect sentence-level reasoning states from multiple reasoning paradigms and tasks;
        \item cluster them into a compact set of global logic states;
        \item train a lightweight logic probe;
        \item compute a routing template for each logic state.
    \end{itemize}
    
    \item \textbf{Online Plug-in Routing}
    \begin{itemize}[leftmargin=1.5em]
        \item extract the current sentence-level reasoning state;
        \item infer the current logic-state distribution;
        \item softly combine cached logic templates;
        \item apply the resulting routing weights to attention heads for the next step.
    \end{itemize}
\end{enumerate}

The full online routing rule is extremely simple:
\[
p_t = \mathrm{Probe}(r_t), \qquad
W_t = \sum_{k=1}^{K} p_t(k)\, M_k,
\]
where
\begin{itemize}[leftmargin=1.5em]
    \item $r_t$ is the current sentence-level reasoning state feature,
    \item $p_t(k)$ is the probability of logic state $k$ at step $t$,
    \item $M_k$ is the cached routing template for logic state $k$.
\end{itemize}

\section{Offline Stage: Global Logic-State Discovery}

\subsection{Step 1: Build a Joint Reasoning-State Pool}

We collect sentence-level reasoning trajectories from diverse reasoning paradigms and tasks. The purpose is not to compare paradigms, but to build a broad and heterogeneous pool of reasoning states from which global logic states can be discovered.

Suppose we have a set of reasoning paradigms
\[
\mathcal{P} = \{\Pi_1,\Pi_2,\dots,\Pi_m\}
\]
and a set of tasks
\[
\mathcal{D} = \{\mathcal{D}_1,\mathcal{D}_2,\dots,\mathcal{D}_n\}.
\]
For each paradigm--task pair, we run the frozen model and extract sentence-level states from generated reasoning trajectories. This yields a pooled collection
\[
\mathcal{R} = \{r_t^{(i)}\}_{i,t},
\]
where each $r_t^{(i)}$ is a sentence-level internal state representation.

This joint pool is important because our hypothesis is that logic is \emph{problem-free}: global logic states are not tied to a specific task or outer reasoning scaffold. Different paradigms may express different subsets of the global logic space, but the space itself should be shared.

\subsection{Step 2: Sentence-Level State Readout}

For each sentence $y_t$, we extract a compact internal state vector
\[
r_t = \Phi(y_t; f_\theta)\in\mathbb{R}^d.
\]

A clean default implementation is to use a late-layer sentence-level residual summary, for example:
\[
r_t = \frac{1}{|y_t|}\sum_{j\in y_t} h_j^{(L)},
\]
where $h_j^{(L)}$ is the token representation at a selected late layer.

Optionally, this state can be augmented by a small number of trajectory-dynamic statistics, but the main method should remain valid even with a simple sentence-level readout.

The point of $r_t$ is not to summarize task semantics alone, but to capture the model's current internal reasoning condition at that step.

\subsection{Step 3: Global Logic Clustering}

We cluster the pooled reasoning states $\mathcal{R}$ into $K$ clusters:
\[
\mathcal{C} = \{C_1,C_2,\dots,C_K\}.
\]

Each cluster corresponds to a global logic state:
\[
z_k \leftrightarrow C_k.
\]

The logic states are intended to be:
\begin{itemize}[leftmargin=1.5em]
    \item \textbf{compact}: small enough to cache and interpret;
    \item \textbf{recurrent}: shared across tasks and paradigms;
    \item \textbf{distinct}: associated with different routing patterns;
    \item \textbf{softly composable}: not forced to be one-hot at test time.
\end{itemize}

We can determine $K$ using a cluster quality criterion such as silhouette score, or by balancing separability and coverage in the pooled reasoning-state space. In the paper, this should be described as a compact global logic basis rather than over-emphasizing any single clustering heuristic.

\subsection{Step 4: Logic Labels as Regimes, Not Rules}

Each cluster defines a \emph{logic regime}. We may assign interpretive names in the analysis section, such as decomposition-like, verification-like, elimination-like, correction-like, or synthesis-like. However, these names should be treated as descriptive summaries rather than hard supervision labels.

Thus, the method does not claim that logic is explicitly given by symbolic rules. Instead, it claims that recurring internal reasoning regimes can be discovered and then used as control units.

\section{Offline Stage: Logic Probe}

\subsection{Pseudo Labels from Clustering}

After clustering, each pooled state $r_t$ receives a pseudo label
\[
\hat{z}_t \in \{1,\dots,K\}.
\]

These pseudo labels define the supervision signal for the logic probe.

\subsection{Probe Definition}

We define a lightweight logic probe:
\[
p_t = \mathrm{softmax}(W r_t + b),
\]
where
\[
p_t \in \Delta^{K-1}.
\]

This gives a probability distribution over the $K$ logic states.

We intentionally use a softmax-based probe rather than independent binary outputs. This encourages competition among logic states and yields sharper online routing. Since our desired routing templates are expected to be sparse and state-specific, a sharper logic distribution is beneficial.

\subsection{Probe Objective}

The probe is trained with a simple cross-entropy loss:
\[
\mathcal{L}_{\text{probe}}
=
-\sum_t \log p_t(\hat{z}_t).
\]

This keeps the controller minimal and avoids introducing additional supervisory components.

\section{Offline Stage: Logic-Specific Routing Templates}

\subsection{Motivation}

The probe alone only tells us which logic state is likely active. We also need to know which internal pathway best supports that logic state. This is the role of the routing template cache.

For each logic state $k$, we want to compute a routing template
\[
M_k \in [0,1]^{L\times H},
\]
where $M_k(\ell,h)$ indicates how strongly head $(\ell,h)$ should participate when the model is in logic state $k$.

\subsection{Head Relevance Under a Logic State}

Let $C_k$ denote the set of reasoning states belonging to logic cluster $k$. For head $(\ell,h)$, define a relevance score under logic state $k$:
\[
I_{\ell,h}^{(k)} = \mathbb{E}_{r_t \in C_k}\big[\mathrm{Rel}_{\ell,h}(r_t)\big].
\]

Here $\mathrm{Rel}_{\ell,h}(r_t)$ measures how functionally aligned head $(\ell,h)$ is with the current logic state.

A natural clean choice is to use a head write-back alignment measure. Let $w_{\ell,h}(t)$ denote the write-back vector of head $(\ell,h)$ at reasoning step $t$. Let $u_k$ denote the centroid direction of logic state $k$ in a shared projected space. Then we can define:
\[
I_{\ell,h}^{(k)}
=
\mathbb{E}_{r_t \in C_k}
\left[
\frac{\langle \tilde{w}_{\ell,h}(t), u_k\rangle^2}{\|\tilde{w}_{\ell,h}(t)\|_2^2+\epsilon}
\right],
\]
where $\tilde{w}_{\ell,h}(t)$ is the projected write-back vector.

This yields a logic-specific head relevance matrix.

\subsection{Template Construction}

We normalize the relevance scores into routing weights:
\[
M_k(\ell,h) = \mathrm{Norm}\!\left(I_{\ell,h}^{(k)}\right),
\]
so that each routing template lies in $[0,1]^{L\times H}$.

These templates are cached offline and reused at inference time.

\subsection{Logic-Insensitive Heads}

A key hypothesis of this work is that many heads are not particularly logic-sensitive. They act as general backbone heads rather than logic-specific controllers.

We quantify logic sensitivity by the variance of each head across logic states:
\[
\sigma_{\ell,h} = \mathrm{Std}\big(\{I_{\ell,h}^{(k)}\}_{k=1}^{K}\big).
\]

If $\sigma_{\ell,h}$ is sufficiently small, then head $(\ell,h)$ is treated as logic-insensitive. Instead of defining a separate whitelist module, we simply absorb such heads into all routing templates by setting
\[
M_k(\ell,h)=1,\qquad \forall k.
\]

This keeps the framework unified and simple:
\begin{itemize}[leftmargin=1.5em]
    \item logic-sensitive heads are selectively modulated;
    \item logic-insensitive heads remain always available.
\end{itemize}

\section{Online Stage: Plug-in Logic Routing}

At inference time, LoTR is attached to a base reasoning paradigm $\Pi$.

\subsection{Step 1: Read Current Sentence-Level State}

After the model completes step $t$, we extract the current sentence-level reasoning state:
\[
r_t = \Phi(y_t; f_\theta).
\]

\subsection{Step 2: Predict Logic-State Distribution}

The logic probe outputs
\[
p_t = \mathrm{softmax}(W r_t + b).
\]

The quantity $p_t(k)$ expresses how strongly the current step aligns with logic state $k$.

\subsection{Step 3: Compose Routing Template}

We compute the final routing weights by soft combination:
\[
W_t = \sum_{k=1}^{K} p_t(k)\, M_k.
\]

This yields a step-specific routing matrix over all layers and heads.

\subsection{Step 4: Modulate Attention Heads}

For each head $(\ell,h)$, let $o_{\ell,h}^{(t)}$ denote its output at the current step. We modulate it as
\[
o_{\ell,h}^{(t)} \leftarrow W_t(\ell,h)\cdot o_{\ell,h}^{(t)}.
\]

No other part of the model is changed. The outer reasoning paradigm remains exactly the same.

\subsection{Interpretation}

This mechanism can be read as follows:
\begin{itemize}[leftmargin=1.5em]
    \item the probe identifies which logic regimes are currently active;
    \item the cached templates encode which heads tend to support those regimes;
    \item the weighted composition suppresses potentially interfering logic pathways while preserving supportive ones.
\end{itemize}

\section{Why the Method Is Clean}

A central strength of LoTR is that it remains structurally simple.

\subsection{Only Two Core Objects}

The entire method is built around two reusable objects:
\begin{enumerate}[leftmargin=1.5em]
    \item a logic probe,
    \item a cache of logic-specific routing templates.
\end{enumerate}

\subsection{Online Computation Is Minimal}

The online routing reduces to:
\[
p_t = \mathrm{Probe}(r_t), \qquad
W_t = \sum_k p_t(k)M_k.
\]
No search, no planner, no extra branch scoring, and no new objective are introduced at inference time.

\subsection{No Redesign of the Base Reasoning Framework}

LoTR does not redefine CoT, SC, SR, or MCTS. It only routes the internal pathway more appropriately at each step.

\subsection{No Handcrafted Logic Rules}

The logic states are discovered from trajectories rather than specified by a symbolic rule set.

\subsection{Finite and Cacheable Control Units}

The clustering step is not for interpretability alone. It is what allows continuous reasoning space to be compressed into a finite cacheable control basis.

\section{Expected Advantages}

We expect LoTR to provide the following benefits.

\subsection{Better Step-Level Reasoning Alignment}

Different reasoning steps require different internal pathway configurations. LoTR adapts pathway participation to the current step rather than relying on a fixed global pattern.

\subsection{Compatibility Across Reasoning Paradigms}

Since it is plug-in, LoTR can be attached to multiple external reasoning frameworks.

\subsection{Suppression of Logic Interference}

If some heads are irrelevant or disruptive under the current logic state, LoTR can down-weight them through the composed routing matrix.

\subsection{Natural Sparsity}

Because many heads are logic-insensitive and many logic distributions are expected to be sharp, the effective routing pattern should be relatively sparse.

\subsection{Interpretability}

The discovered logic states and their corresponding routing templates provide a natural analysis interface:
\begin{itemize}[leftmargin=1.5em]
    \item which logic regimes recur across paradigms?
    \item which heads are logic-sensitive?
    \item which heads are shared backbones?
    \item which routing patterns help which reasoning steps?
\end{itemize}

\section{Main Experimental Plan}

\subsection{Part I: Plug-in Effectiveness Across Reasoning Paradigms}

\textbf{Goal.} Show that LoTR works as a general plug-in controller.

\textbf{Base paradigms.}
\begin{itemize}[leftmargin=1.5em]
    \item Vanilla
    \item CoT
    \item Self-Consistency (SC)
    \item Self-Refine (SR)
    \item Plan-and-Solve
    \item MCTS-style reasoning
\end{itemize}

\textbf{Comparison.}
For each paradigm, compare:
\begin{itemize}[leftmargin=1.5em]
    \item base paradigm,
    \item base paradigm + LoTR.
\end{itemize}

\textbf{Metrics.}
\begin{itemize}[leftmargin=1.5em]
    \item accuracy / exact match,
    \item pass@1 or recall where appropriate,
    \item average token usage,
    \item latency,
    \item cost-adjusted performance,
    \item robustness under more difficult settings.
\end{itemize}

\textbf{Target conclusion.}
LoTR should improve multiple reasoning paradigms without requiring paradigm-specific redesign.

\subsection{Part II: Logic-State Discovery Analysis}

\textbf{Goal.} Show that the discovered clusters are meaningful as recurring reasoning regimes.

\textbf{Analyses.}
\begin{itemize}[leftmargin=1.5em]
    \item visualization of the logic-state space,
    \item distribution of different paradigms over logic states,
    \item distribution of tasks over logic states,
    \item qualitative interpretation of each cluster.
\end{itemize}

\textbf{Target conclusion.}
The discovered logic states are neither arbitrary nor paradigm-specific; they recur across diverse trajectories.

\subsection{Part III: Routing Specificity}

\textbf{Goal.} Show that logic-specific templates matter.

\textbf{Key experiments.}
\begin{itemize}[leftmargin=1.5em]
    \item matched template vs random template,
    \item matched template vs average template,
    \item off-diagonal swap matrix: use template $M_j$ on steps classified as logic $i$,
    \item shuffled cluster labels.
\end{itemize}

\textbf{Target conclusion.}
The gain does not come from generic reweighting alone, but from matching the current logic state to the right routing template.

\subsection{Part IV: Not Just Static Pruning or Surface Features}

\textbf{Goal.} Distinguish LoTR from simpler alternatives.

\textbf{Baselines.}
\begin{itemize}[leftmargin=1.5em]
    \item static head pruning,
    \item average head importance,
    \item activation-only routing,
    \item routing by sentence length,
    \item routing by step index,
    \item routing by task ID,
    \item routing by paradigm ID.
\end{itemize}

\textbf{Target conclusion.}
LoTR should outperform baselines that only exploit static salience or surface metadata.

\subsection{Part V: Transferability}

\textbf{Goal.} Support the claim that logic is problem-free.

\textbf{Experiments.}
\begin{itemize}[leftmargin=1.5em]
    \item train logic states on a subset of paradigms, test on unseen paradigms,
    \item train on a subset of tasks, test on unseen tasks,
    \item optionally test across nearby backbones.
\end{itemize}

\textbf{Target conclusion.}
LoTR captures reusable logic regimes rather than overfitting to one task or one scaffold.

\subsection{Part VI: Sparsity and Efficiency}

\textbf{Goal.} Validate the intended routing structure.

\textbf{Metrics.}
\begin{itemize}[leftmargin=1.5em]
    \item distribution sharpness of $p_t$,
    \item proportion of logic-sensitive vs logic-insensitive heads,
    \item overlap between routing templates,
    \item number of strongly suppressed heads per step,
    \item online overhead.
\end{itemize}

\textbf{Target conclusion.}
The method is lightweight, sparse in practice, and structurally consistent with its design hypothesis.

\section{Critical Ablations}

The following ablations are particularly important.

\subsection{No Clustering}

Replace the discrete logic basis with a fully continuous controller.

\textbf{Purpose.} Show that a compact finite logic basis is useful for cacheability, interpretability, and stable routing.

\subsection{No Softmax Competition}

Replace softmax with non-competitive outputs.

\textbf{Purpose.} Show that sharper logic-state competition improves routing specificity.

\subsection{No Logic-Insensitive Absorption}

Force all heads to be logic-dependent.

\textbf{Purpose.} Show that preserving invariant backbone heads is beneficial.

\subsection{Average Template Only}

Use a single global template for all steps.

\textbf{Purpose.} Show that dynamic logic-conditioned routing is the key, not just global reweighting.

\subsection{Random or Surface-Level Clustering}

Cluster using random assignments or obvious surface features.

\textbf{Purpose.} Show that the quality of the logic partition is central.

\section{Potential Risks and How to Address Them}

\subsection{Risk 1: Clusters May Reflect Surface Style Rather Than Logic}

This is a legitimate concern. The intended response is not to claim perfect semantic logic labeling, but to show that the discovered regimes are functionally meaningful through:
\begin{itemize}[leftmargin=1.5em]
    \item recurrence across paradigms,
    \item recurrence across tasks,
    \item matched-routing superiority,
    \item superiority over surface-feature baselines.
\end{itemize}

\subsection{Risk 2: Gains May Be Interpreted as Generic Dynamic Pruning}

This should be addressed by demonstrating:
\begin{itemize}[leftmargin=1.5em]
    \item LoTR operates at sentence-level logic adaptation,
    \item matched routing is better than random or average routing,
    \item static pruning baselines do not recover the same gains,
    \item routing is driven by internal state rather than fixed scenario metadata.
\end{itemize}

\subsection{Risk 3: Pseudo Labels from Clustering May Be Noisy}

This is true in principle, but the method does not depend on exact semantic labels. It depends on whether the partition yields reusable and functionally distinct routing templates. This is testable through the routing-specificity experiments.

\subsection{Risk 4: The Controller May Add Complexity Without Clear Benefit}

This is why the paper should keep the online controller maximally simple and push all complexity into offline compilation. The clean formula
\[
W_t = \sum_k p_t(k)M_k
\]
must remain the conceptual center of the method.

\section{Main Contribution Summary}

If successful, this paper would contribute the following:

\begin{enumerate}[leftmargin=1.5em]
    \item \textbf{A new plug-in reasoning controller.}  
    LoTR introduces step-level logic-conditioned routing without replacing existing reasoning paradigms.
    
    \item \textbf{A compact global logic basis.}  
    It proposes to discover a finite set of recurring logic states from heterogeneous reasoning trajectories.
    
    \item \textbf{A simple logic--pathway coupling mechanism.}  
    Each logic state is associated with a reusable routing template over attention heads.
    
    \item \textbf{A new view of reasoning enhancement.}  
    The work shifts emphasis from external reasoning scaffolds alone to matching the model's current internal condition with a more suitable inference pathway.
\end{enumerate}

\section{Concise One-Paragraph Abstract Draft}

We propose \textbf{Logic-of-Thought Routing (LoTR)}, a plug-in controller for inference-time reasoning that dynamically adjusts the model's internal pathway according to its current step-level logic state. Existing reasoning methods often rely on externally injected programs, such as prompting strategies, refinement loops, or search procedures, but they do not explicitly adapt the model's internal computation to the current reasoning step. LoTR addresses this gap by discovering a compact set of global logic states from heterogeneous reasoning trajectories, training a lightweight probe to predict the current logic-state distribution from internal information transfer, and softly combining cached logic-specific routing templates to modulate attention heads. The method is simple, reusable, and compatible with diverse reasoning paradigms including CoT, self-consistency, self-refinement, and search-based reasoning. We hypothesize that different logic states are coupled with different sparse attention-head participation patterns, and that routing heads according to the current logic state can suppress interfering logic pathways while preserving supportive ones. If validated, LoTR would introduce a clean new interface for plug-in reasoning: not changing the outer reasoning scaffold, but improving reasoning by matching each step to a more suitable internal pathway.

\section{Conclusion}

LoTR is designed to be a clean and modular reasoning proposal. It does not attempt to hand-code logic or replace existing reasoning methods. Instead, it introduces a minimal internal routing layer that reads the model's current reasoning state, maps it to a compact global logic space, and activates a more suitable pathway for the next step. The main promise of this work is not complexity, but simplicity: a single lightweight probe, a finite cache of routing templates, and a stepwise soft routing rule that can be attached to many reasoning paradigms. If the empirical evidence supports the central hypotheses, LoTR could provide a new and elegant direction for plug-in reasoning.

\end{document}